\chapter{Procedimientos normales}
\label{cap:normales}

\opening{La rutina bien hecha es la que deja energía libre para lo inesperado.}

\section{Velocidades recomendadas}
\fuente{POH, sección 4 (aprobada por la FAA)}

\begin{checklist}
\chk{Despegue y ascenso}{60 KIAS}
\chk{Máxima razón de ascenso (\textsc{v}\textsubscript{Y})}{53 KIAS}
\chk{Máximo alcance}{83 KIAS\textsuperscript{*}}
\chk{Turbulencia significativa}{60 a 70 KIAS\textsuperscript{*}}
\chk{Aproximación para aterrizar}{60 KIAS}
\chk{Autorrotación}{60 a 70 KIAS\textsuperscript{*}}
\end{checklist}

{\footnotesize\textsuperscript{*} Ciertas condiciones pueden requerir una velocidad menor. Vea los letreros de \textsc{v}\textsubscript{NE} en el capítulo~\ref{cap:limitaciones}.}

\section{Chequeos diarios o de prevuelo}

Antes del primer vuelo del día, y de cada vuelo, debe completarse la inspección diaria o de prevuelo del fabricante. Por su extensión, la lista completa, desarrollada ítem por ítem, constituye la Parte~III de este libro (capítulos~\ref{cap:prevuelo-guia} y~\ref{cap:prevuelo-lista}). Sus procedimientos generales son:

\begin{itemize}
\item Retire las ruedas de maniobra en tierra y todas las fundas y amarras.
\item Retire incluso pequeñas acumulaciones de escarcha, hielo o nieve, especialmente de las palas del rotor.
\item Revise los registros de mantenimiento para verificar que el helicóptero esté aeronavegable.
\item Revise el estado general y verifique que no haya daños visibles, pérdidas de fluidos ni desgaste anormal.
\item Verifique que no haya \emph{fretting} (desgaste por micromovimiento) en remaches y uniones. El \emph{fretting} del aluminio produce un polvo negro fino; el del acero, un residuo café rojizo o negro.
\item Verifique que los Telatemps no muestren aumentos de temperatura no atribuibles a un cambio en las condiciones de operación (los mecánicos trazan una línea de referencia a la derecha del cuadro de temperatura más alta que se haya oscurecido en operación).
\item Verifique que las marcas de torque de los sujetadores críticos no estén rotas ni falten.
\end{itemize}

\section{Antes de arrancar el motor}

\begin{checklist}
\chk{Cinturones de seguridad}{Abrochados}
\chk{Válvula de corte de combustible}{ON}
\chk{Fricción de cíclico/colectivo}{OFF}
\chk{Cíclico, colectivo y pedales}{Recorrido completo libre}
\chk{Acelerador}{Recorrido completo libre}
\chk{Colectivo}{Completamente abajo, fricción ON}
\chk{Cíclico}{Neutro, fricción ON}
\chk{Pedales}{Neutros}
\chk{Freno rotor}{Desaplicado}
\chk{Disyuntores}{Adentro}
\chk{Calefacción del carburador}{OFF}
\chk{Mezcla}{Rica a fondo}
\chk{Protector de la mezcla\textsuperscript{*}}{Instalado}
\chk{Cebador (si está instalado)}{Abajo y trabado}
\chk{Luces de aterrizaje}{OFF}
\chk{Interruptor de aviónica (si está instalado)}{OFF}
\chk{Embrague}{Desacoplado}
\chk{Altímetro}{Ajustado}
\chk{Interruptor del gobernador}{ON}
\end{checklist}

{\footnotesize\textsuperscript{*} El protector de la mezcla no se usa en los helicópteros con control de mezcla tipo vernier en la cara de la consola.}

\section{Arranque del motor y prueba}

\begin{checklist}
\chk{Giros de acelerador para cebar}{Según se requiera}
\chk{Acelerador}{Cerrado}
\chk{Interruptores de batería y estrobo}{ON}
\chk{Área}{Despejada}
\chk{Interruptor de encendido}{START, luego BOTH}
\chk{Luz \textsf{STARTER ON}}{Apagada}
\chk{RPM del motor}{Ajustar 50 a 60\,\%}
\chk{Interruptor del embrague}{Acoplado (\textsf{ENGAGE})}
\chk{Palas girando}{En menos de 5 s}
\chk{Interruptor del alternador}{ON}
\chk{Presión de aceite dentro de 30 s}{25 psi mínimo}
\chk{Aviónica, audífonos}{ON}
\chk{Alerta sonora (si está equipada)}{Probar}
\chk{Esperar que se apague la luz del embrague}{Disyuntores adentro}
\chk{RPM de calentamiento}{70 a 75\,\%}
\chk{Instrumentos de motor}{En verde}
\chk{Caída de magnetos a 75\,\% RPM}{7\,\% máx. en 2 s}
\chk{Calefacción del carburador}{Subida/caída de CAT; ajustar según se requiera}
\chk{Prueba del embrague sprag}{Agujas se separan}
\chk{Puertas (si están instaladas)}{Cerradas y trabadas}
\chk{Gráfico de límite de MAP}{Revisar}
\chk{Fricción de cíclico/colectivo}{OFF}
\chk{Gobernador ON, aumentar acelerador}{RPM 102--104\,\%}
\chk{Luces de advertencia}{Apagadas}
\chk{Subir levemente el colectivo, reducir RPM}{Bocina/luz a 97\,\%}
\end{checklist}

\begin{precaucion}
En los helicópteros que entregan la bocina de baja RPM a través del sistema de audio, cada piloto necesita un audífono para escuchar la bocina.
\end{precaucion}

\begin{precaucion}
Evite la operación continua con una RPM del rotor de 60 a 70\,\% para minimizar la resonancia del cono de cola.
\end{precaucion}

\begin{precaucion}
En superficies resbaladizas, esté preparado para contrarrestar con pedal izquierdo la rotación de nariz a la derecha cuando el gobernador aumente la RPM.
\end{precaucion}

\section{Despegue}

\begin{pasos}
\item Verifique puertas trabadas, gobernador ON y RPM estabilizada entre 102 y 104\,\%.
\item Despeje el área. Suba lentamente el colectivo hasta que el helicóptero quede liviano sobre los patines. Reposicione el cíclico según se requiera para el equilibrio y luego eleve suavemente el helicóptero a vuelo estacionario.
\item Revise que los instrumentos estén en verde y ajuste la calefacción del carburador si es necesario.
\item Baje la nariz y acelere a la velocidad de ascenso siguiendo el perfil del diagrama altura--velocidad (capítulo~\ref{cap:performance}). Si la RPM cae por debajo de 102\,\%, baje el colectivo.
\end{pasos}

\section{Crucero}

\begin{pasos}
\item Ajuste la calefacción del carburador si es necesario (sección~\ref{sec:carb-heat}).
\item Verifique la RPM cerca del extremo superior del arco verde.
\item Ajuste con el colectivo la presión de admisión deseada. Respete los límites de MAP y de velocidad.
\item Tire del pomo de compensación \textsf{RT TRIM}.
\item Verifique instrumentos en verde y luces de advertencia apagadas.
\end{pasos}

\begin{precaucion}
En turbulencia, reduzca la potencia y use una velocidad de crucero menor que la normal. Si la turbulencia es significativa o se vuelve incómoda para el piloto, use 60 a 70~KIAS.
\end{precaucion}

\begin{precaucion}
Tenga extremo cuidado de no tirar nunca inadvertidamente del control de mezcla: el motor se detendrá.
\end{precaucion}

\begin{precaucion}
No se recomienda empobrecer la mezcla en vuelo. El motor puede detenerse, ya que no hay hélice que lo mantenga girando si se empobrece en exceso.
\end{precaucion}

\section{Operación sin puertas}

Evite quitar la puerta izquierda para proteger el rotor de cola de objetos sueltos. Si la puerta izquierda debe quitarse, advierta al pasajero que asegure los objetos sueltos y que mantenga la cabeza y los brazos dentro de la cabina para evitar la corriente de aire de alta velocidad.

\section{Autorrotación de práctica con recuperación de potencia}

\begin{pasos}
\item Ajuste la calefacción del carburador si es necesario.
\item Baje el colectivo hasta el tope inferior y reduzca el acelerador lo necesario para separar las agujas del tacómetro.
\end{pasos}

\begin{precaucion}
Para evitar una detención inadvertida del motor, no corte el acelerador de golpe para simular una falla de potencia. Siempre ciérrelo suavemente. Recupere de inmediato si el motor funciona irregularmente o si la RPM del motor sigue cayendo.
\end{precaucion}

\begin{nota}
El gobernador está inactivo por debajo de 80\,\% de RPM del motor, cualquiera sea la posición de su interruptor.
\end{nota}

\begin{nota}
Al entrar en autorrotación sobre 4000~ft, reduzca levemente el acelerador antes de bajar el colectivo para evitar una sobrevelocidad del motor.
\end{nota}

\begin{pasos}
\setcounter{enumi}{2}
\item Ajuste el colectivo para mantener la RPM del rotor dentro de límites y el acelerador para mantener las agujas separadas.
\item Mantenga la velocidad entre 60 y 70~KIAS.
\item A unos 40~ft AGL, inicie la recogida con el cíclico para reducir la razón de descenso y la velocidad de avance.
\item A unos 8~ft AGL, aplique cíclico adelante para nivelar el helicóptero y suba el colectivo para controlar el descenso. Agregue acelerador si es necesario para mantener la RPM en el arco verde.
\end{pasos}

\begin{precaucion}
Las fallas de motor simuladas requieren bajar el colectivo con prontitud para evitar una RPM del rotor peligrosamente baja. Podría producirse una pérdida catastrófica del rotor si la RPM cae alguna vez por debajo de 80\,\% más 1\,\% por cada 1000~ft de altitud.
\end{precaucion}

\section{Autorrotación de práctica con contacto con el suelo}

Si se requieren autorrotaciones con contacto con el suelo para fines de demostración, efectúelas igual que las de recuperación de potencia, excepto:

\begin{itemize}
\item Antes de la recogida, gire el acelerador hasta cerrarlo contra el resorte de sobrerrecorrido y manténgalo contra el tope duro hasta completar la autorrotación. (Esto impide que el correlador del acelerador agregue potencia al subir el colectivo.)
\item Toque siempre el suelo con los patines nivelados y la nariz recta hacia adelante.
\end{itemize}

\begin{precaucion}
El R22 tiene un sistema rotor liviano y de baja inercia. La mayor parte de la energía necesaria para una autorrotación está almacenada en la cantidad de movimiento de avance del helicóptero, no en el rotor. Por lo tanto, se requiere una recogida con el cíclico bien sincronizada y la RPM del rotor debe mantenerse en verde hasta justo antes del contacto con el suelo.
\end{precaucion}

\begin{nota}
Haga inspeccionar con frecuencia las zapatas de los patines al practicar autorrotaciones con contacto con el suelo: pueden desgastarse rápidamente.
\end{nota}

\section{Uso de la calefacción del carburador}
\label{sec:carb-heat}

El hielo en el carburador puede formarse en una amplia gama de condiciones atmosféricas, pero es más probable cuando la OAT está entre $-4$~°C y 30~°C (25~°F y 86~°F) y la diferencia entre la OAT y el punto de rocío es menor que 15~°C (27~°F). Cuando se sospechen condiciones propicias para el hielo en el carburador, use la calefacción así:

\begin{description}
\item[Durante la prueba de motor:] use calefacción total (el aire está filtrado) durante el calentamiento para precalentar el sistema de inducción.
\item[Durante el despegue, el ascenso y el crucero:] use la calefacción necesaria para mantener la indicación de CAT fuera del arco amarillo.
\item[Durante el descenso y la autorrotación:] con potencias por debajo de 18~inHg de MAP, aplique calefacción total independientemente de la indicación de CAT. El indicador de CAT no indica correctamente la temperatura del carburador por debajo de 18~inHg de MAP.
\end{description}

\begin{precaucion}
El piloto puede no darse cuenta de la formación de hielo en el carburador, porque el gobernador abre automáticamente el acelerador y mantiene constantes la presión de admisión y la RPM. Por lo tanto, el piloto debe aplicar calefacción al carburador según sea necesario siempre que se sospechen condiciones de engelamiento.
\end{precaucion}

\section{Uso del asistente de calefacción del carburador}

Los R22 con motor O-360 tienen instalado un asistente de calefacción del carburador (\emph{carb heat assist}). Este dispositivo correlaciona la aplicación de calefacción con los cambios de colectivo para reducir la carga de trabajo del piloto: bajar el colectivo agrega calefacción mecánicamente y subirlo la reduce. Un embrague de fricción permite al piloto sobreponerse al sistema y aumentar o disminuir la calefacción según sea necesario.

El pomo de control tiene un pestillo para trabar la calefacción en OFF. El pomo debe dejarse destrabado salvo que sea obvio que las condiciones no son propicias para el hielo en el carburador. Aplique calefacción según sea necesario si existe la posibilidad de hielo. Vigile el indicador de CAT y reajuste según sea necesario después de elevarse a estacionario o de cualquier cambio de potencia.

\section{Descenso, aproximación y aterrizaje}

\begin{pasos}
\item Reduzca la potencia con el colectivo según desee. Ajuste la calefacción del carburador según sea necesario. Respete los límites de velocidad.
\end{pasos}

\begin{precaucion}
No inicie un descenso con cíclico adelante: puede producir una condición de baja G. Inicie siempre el descenso bajando el colectivo.
\end{precaucion}

\begin{pasos}
\setcounter{enumi}{1}
\item Haga la aproximación final contra el viento, con la menor razón de descenso practicable y una velocidad inicial de 60~nudos.
\item Reduzca suavemente la velocidad y la altura hasta el estacionario. (Asegúrese de que la razón de descenso sea menor que 300~ft/min antes de reducir la velocidad por debajo de 30~KIAS.)
\item Desde el estacionario, baje gradualmente el colectivo hasta el contacto con el suelo.
\item Después del contacto inicial, baje el colectivo hasta su posición completamente abajo.
\end{pasos}

\begin{precaucion}
Al aterrizar en pendiente, devuelva el cíclico a neutro antes de reducir la RPM del rotor.
\end{precaucion}

\begin{precaucion}
Nunca deje los controles de vuelo desatendidos con el motor en marcha.
\end{precaucion}

\begin{precaucion}
Mantenga el acelerador cerrado si un pasajero sube o baja con el motor en marcha y el colectivo izquierdo instalado.
\end{precaucion}

\section{Detención del motor}

\begin{checklist}
\chk{Colectivo abajo, RPM 70--75\,\%}{Fricción ON}
\chk{Cíclico y pedales neutros}{Fricción ON}
\chk{Caída de CHT}{Acelerador cerrado}
\chk{Interruptor del embrague}{Desacoplar}
\chk{Esperar 30 s}{Mezcla OFF}
\chk{Protector de la mezcla}{De vuelta sobre la mezcla}
\chk{Esperar 30 s}{Aplicar freno rotor}
\chk{Luz del embrague}{Se apaga}
\chk{Interruptores de aviónica, alternador, batería y encendido}{OFF}
\end{checklist}

\begin{nota}
Si la temperatura ambiente supera los 100~°F (38~°C), enfríe el motor a 70--75\,\% de RPM durante al menos un minuto antes de reducir a ralentí.
\end{nota}

\begin{nota}
Durante el ralentí y después de detener el motor, el piloto debe descubrirse un oído y escuchar cualquier ruido inusual que pueda indicar la falla inminente de un rodamiento u otro componente.
\end{nota}

\begin{precaucion}
No frene el rotor subiendo el colectivo durante la detención. Las palas pueden aletear y golpear el cono de cola.
\end{precaucion}

\section{Atenuación de ruido}

Para mejorar la calidad del entorno y evitar ordenanzas excesivamente restrictivas contra los helicópteros, es imprescindible que cada piloto minimice las molestias por ruido al público. Cuando sea posible, emplee las siguientes técnicas:

\begin{pasos}
\item Evite volar sobre aglomeraciones de personas al aire libre. Cuando no pueda evitarse, vuele tan alto como sea práctico, preferiblemente sobre 2000~ft AGL.
\item Evite el golpeteo de palas (\emph{blade slap}). Suele producirse en descensos poco pronunciados a alta velocidad, especialmente en virajes, y se evita con descensos más lentos y empinados. Con la puerta derecha removida, el piloto puede identificar fácilmente las condiciones de vuelo que lo producen y desarrollar técnicas para eliminarlo o reducirlo.
\item Al salir de un lugar de aterrizaje o aproximarse a él, evite el vuelo prolongado sobre áreas sensibles al ruido. Vuele siempre sobre 500~ft AGL y preferiblemente sobre 1000~ft AGL.
\item El ruido repetitivo irrita mucho más que un evento aislado. Si debe sobrevolar la misma zona más de una vez, varíe la trayectoria para no pasar sobre los mismos edificios.
\item Al sobrevolar áreas pobladas, mire hacia adelante y elija la ruta menos sensible al ruido.
\end{pasos}

\begin{nota}
Estos procedimientos no aplican cuando contradigan autorizaciones o instrucciones del control de tránsito aéreo, o cuando a juicio del piloto produzcan una trayectoria insegura.
\end{nota}

\section{Información según la AD 95-26-04 de la FAA}

Hasta que la FAA complete su investigación de las condiciones y características que conducen a los accidentes por contacto de las palas del rotor principal con el fuselaje, y se identifiquen cambios de diseño y limitaciones de operación correctivas, se insta encarecidamente a los pilotos de R22 a familiarizarse con la siguiente información y a cumplir los procedimientos recomendados.

\textbf{Pérdida (\emph{stall}) del rotor principal.} Muchos factores pueden contribuir y el piloto debe conocerlos. Cualquier condición de vuelo que cree un ángulo de ataque excesivo en las palas puede producir una pérdida. La baja RPM del rotor, las maniobras agresivas, un ángulo de colectivo alto (frecuentemente por alta altitud de densidad, exceso de paso ---exceder la potencia disponible--- en ascenso o alta velocidad de avance) y la respuesta lenta a la bocina y la luz de baja RPM pueden producir la pérdida del rotor principal. El efecto de estas condiciones puede amplificarse en turbulencia. La pérdida del rotor principal puede terminar en el contacto entre el rotor y el fuselaje. Las Safety Notices SN-10, SN-15, SN-20, SN-24, SN-27 y SN-29 entregan información adicional.

\textbf{Golpe de mástil (\emph{mast bumping}).} Puede ocurrir en un rotor balancín cuando hay un aleteo excesivo del rotor principal por baja G (factor de carga menor que 1,0) o por una acción brusca de los controles. Una condición de baja G puede resultar de un empuje brusco del cíclico en vuelo de avance. La alta velocidad, la turbulencia y el derrape excesivo pueden acentuar los efectos adversos de estos movimientos. El aleteo excesivo hace que el cubo del rotor golpee el mástil, con la posterior separación del sistema rotor del helicóptero.

Para evitar estas condiciones, se recomienda encarecidamente:

\begin{pasos}
\item Mantener velocidades de crucero entre 60~KIAS y menos de 0,9\,\textsc{v}\textsubscript{NE}, pero no menos de 57~KIAS.
\item Usar la RPM máxima ``con potencia'' en todo momento durante el vuelo con potencia.
\item Evitar el derrape en vuelo. Mantener el vuelo compensado en todo momento.
\item Evitar acciones grandes y rápidas de cíclico adelante en vuelo de avance, y acciones bruscas de los controles en turbulencia.
\end{pasos}
